\chapter{Experimental Results}
\label{chap:results}

Every ordinary panel in this chapter begins at a true MI difference of zero.
That point measures false-positive rejection under \(H_0\). Points to its
right measure detection under \(H_1\). The horizontal coordinate is a
\emph{constructed population} difference in nats; random count tables create
the vertical variation. The blue solid circles are Normal Wald and the
magenta dashed squares are Expanded Welch. The dotted horizontal line marks
the nominal 0.05 rejection rate. Each plotted rate divides the number of valid rejections
by all 20,000 replicates. A hollow marker indicates a valid-result rate below
0.90, not an accepted test. The figure source manifest gives the exact saved
configuration identifiers for every panel.

\section{Main landscape: shape, margins and sample size}

The first question is broad: do the methods behave differently as the table
grows, margins become uneven and samples become larger?
Figures~\ref{fig:main-2x2-uniform}--\ref{fig:main-8x8-different_skew} provide a
systematic view, not a selection of favourable examples. Each of the
four shapes is shown under all three margin profiles. Within every figure,
the columns are the same three equal sample sizes. The full vertical scale
preserves comparability of detection; the lower row enlarges the region in
which calibration differences near 0.05 can be seen. The complete nine-size
ladder, including \(n=2,5,20,50,250,500\), remains in the companion atlas.

\newcommand{\mainlandscapefigure}[5]{%
\begin{figure}[p]
\centering
\makebox[\textwidth][c]{\includegraphics[width=1.18\textwidth]{main_#1_#3.pdf}}
\caption[Main landscape: #2 tables, #5 margins]{#2 tables, #4. Columns: \(n_P=n_Q=10,100,1000\), respectively.
For each population the first two rows and columns receive the signed changes
in Equation~\eqref{eq:first-change-block}; \(I(P)=0.02\) and
\(I(Q)=0.02+|\Delta|\). The horizontal grid is
\(|\Delta|=0,0.001,0.002,0.005,0.01,0.02\) nats. The upper panels have a
shared 0--1 rejection axis and the lower panels show the same data at 0--0.15.
At zero the rate is false positives; elsewhere it is detection. Each point
uses 20,000 paired samples. Hollow markers denote validity below 90\%.}
\label{fig:main-#1-#3}
\end{figure}%
}

\mainlandscapefigure{2x2}{\(2\times2\)}{uniform}{both \(P\) and \(Q\) have uniform row and column margins \((1/2,1/2)\)}{uniform}
\mainlandscapefigure{2x2}{\(2\times2\)}{same_skew}{both \(P\) and \(Q\) have row and column margins \(m_2(0.8)=(0.8,0.2)\)}{same skew}
\mainlandscapefigure{2x2}{\(2\times2\)}{different_skew}{\(P\) has row and column margins \(m_2(0.7)=(0.7,0.3)\), while \(Q\) has \(m_2(0.8)=(0.8,0.2)\)}{different skew}

\mainlandscapefigure{3x3}{\(3\times3\)}{uniform}{both \(P\) and \(Q\) have uniform row and column margins \((1/3,1/3,1/3)\)}{uniform}
\mainlandscapefigure{3x3}{\(3\times3\)}{same_skew}{both \(P\) and \(Q\) have row and column margins \(m_3(0.8)=(0.8,0.1,0.1)\)}{same skew}
\mainlandscapefigure{3x3}{\(3\times3\)}{different_skew}{\(P\) has row and column margins \(m_3(0.7)=(0.7,0.15,0.15)\), while \(Q\) has \(m_3(0.8)=(0.8,0.1,0.1)\)}{different skew}

\mainlandscapefigure{5x5}{\(5\times5\)}{uniform}{both \(P\) and \(Q\) have uniform row and column margins \((1/5,\ldots,1/5)\)}{uniform}
\mainlandscapefigure{5x5}{\(5\times5\)}{same_skew}{both \(P\) and \(Q\) have row and column margins \(m_5(0.8)=(0.8,0.05,0.05,0.05,0.05)\)}{same skew}
\mainlandscapefigure{5x5}{\(5\times5\)}{different_skew}{\(P\) has row and column margins \(m_5(0.7)=(0.7,0.075,0.075,0.075,0.075)\), while \(Q\) has \(m_5(0.8)=(0.8,0.05,0.05,0.05,0.05)\)}{different skew}

\mainlandscapefigure{8x8}{\(8\times8\)}{uniform}{both \(P\) and \(Q\) have uniform row and column margins \((1/8,\ldots,1/8)\)}{uniform}
\mainlandscapefigure{8x8}{\(8\times8\)}{same_skew}{both \(P\) and \(Q\) have row and column margins \(m_8(0.8)=(0.8,0.2/7,\ldots,0.2/7)\)}{same skew}
\mainlandscapefigure{8x8}{\(8\times8\)}{different_skew}{\(P\) has row and column margins \(m_8(0.7)=(0.7,0.3/7,\ldots,0.3/7)\), while \(Q\) has \(m_8(0.8)=(0.8,0.2/7,\ldots,0.2/7)\)}{different skew}

The landscape does not yield a uniformly preferable method. At \(n=5\),
Expanded Welch rejects none of the \(\MainFiveNullCount\) main null regimes. This includes the
uniform \(3\times3\), \(5\times5\) and \(8\times8\) regimes, where more than
90\% of calculations are valid, so the zero rates cannot be attributed only
to invalid outputs. Among the \(\MainTenTwentyNullCount\) main null regimes with \(n=10\) or 20,
Expanded Welch is numerically closer to 0.05 in
\(\MainTenTwentyExpandedCloserCount\) and Wald in
\(\MainTenTwentyWaldCloserCount\). This is a
descriptive count, not a test of independent improvements. The trend reverses
in the \(\MainLateNullCount\) main null settings at \(n=500\) or 1,000: Wald
is closer to 0.05 in all \(\MainLateWaldCloserCount\). At \(n=1{,}000\), Wald
null rates across the main settings lie between
\(\MainThousandWaldMinRate\) and \(\MainThousandWaldMaxRate\); Expanded
Welch lies between \(\MainThousandExpandedMinRate\) and
\(\MainThousandExpandedMaxRate\). Neither is uniformly at the nominal
level.

The interpretation is directional. Expanded Welch uses a heavier-tailed
reference and therefore lowers rejection; it does not diagnose whether Wald
started above or below 0.05. The same adjustment can consequently improve a
liberal regime and worsen a conservative one. Closeness to 0.05 must be read
regime by regime and not treated as an inherent advantage of either method.

Two fully valid null points make the contrast concrete. With uniform
\(8\times8\) populations, \(n_P=n_Q=20\), and equal true MI of 0.02,
Wald rejects \(\EightUniformTwentyWaldRate\) of samples and Expanded Welch
\(\EightUniformTwentyExpandedRate\). The latter is closer to 0.05. With
uniform \(2\times2\) populations at \(n_P=n_Q=1{,}000\), the respective
rates are \(\TwoUniformThousandWaldRate\) and
\(\TwoUniformThousandExpandedRate\); the correction moves an already
conservative procedure farther below the target. Both methods are valid on
all replicates in these examples, so the change is attributable to their
different references, not missing outputs.

These two examples isolate the central trade-off. In the \(8\times8\) case,
the lower rejection rate corrects part of Wald's excess. In the \(2\times2\)
case, the same downward movement is harmful because Wald was already
conservative. Expanded Welch is therefore a one-direction adjustment, not an
adaptive calibration procedure.

\begin{table}[htbp]
\centering
\small
\caption{Two fully valid main-grid null regimes. Brackets beside each rate
are pointwise Wilson 95\% intervals; the final bracket is the approximate
paired 95\% interval for Wald minus Expanded Welch on the same samples.
Source configurations are \texttt{config\_313a5e5c05280f9e} and
\texttt{config\_700945d218841f2e}, respectively.}
\label{tab:null-paired-examples}
\begin{tabularx}{\textwidth}{@{}l>{\centering\arraybackslash}X>{\centering\arraybackslash}X>{\centering\arraybackslash}X@{}}
\toprule
Regime & Wald & Expanded & Wald minus Expanded\\
 & Rate [95\% interval] & Rate [95\% interval] & Rate [95\% interval]\\
\midrule
\(8\times8\), uniform, \(n=20\) & 0.0687 [0.0653, 0.0723] & 0.0570 [0.0538, 0.0602] & 0.0118 [0.0103, 0.0132]\\
\(2\times2\), uniform, \(n=1000\) & 0.0445 [0.0417, 0.0474] & 0.0387 [0.0361, 0.0414] & 0.0059 [0.0048, 0.0069]\\
\bottomrule
\end{tabularx}
\end{table}

The lower sample-size figures also show an important limit to interpreting
small rejection rates. In the \(8\times8\) different-skew alternative with
\(n_P=n_Q=5\), \(I(P)=0.02\) and \(I(Q)=0.021\), Wald rejects
\(\EightSparseFiveAlternativeWaldRate\) and Expanded Welch rejects
\(\EightSparseFiveAlternativeExpandedRate\), but their valid-result rates
are only \(\EightSparseFiveAlternativeWaldValid\) and
\(\EightSparseFiveAlternativeExpandedValid\). These are probabilities of a
\emph{valid rejection}, not a clean comparison of power at matched achieved
false-positive rates. The corresponding null and validity records must be
examined before calling the gap a power advantage.

\section{Baseline MI and the location of changed cells}

The main grid fixes \(I(P)=0.02\). Does starting closer to independence
change the conclusion? Figure~\ref{fig:baseline-focus} holds shape, margins,
sample size and changed cells fixed while varying only the common starting
MI\@. The zero-difference points are especially informative: for
\(I_0=0.0001,0.001,0.02\), Wald's rates are
\(\BaselineTinyWaldRate\), \(\BaselineSmallWaldRate\) and
\(\BaselineRegularWaldRate\), while Expanded Welch's are
\(\BaselineTinyExpandedRate\), \(\BaselineSmallExpandedRate\) and
\(\BaselineRegularExpandedRate\). All the plotted null
calculations in this focused set are valid. Replacing the normal reference
with a heavier-tailed Student reference does not address the already-low null
rejection in these settings.

Starting closer to independence therefore does not create a general case for
Expanded Welch. Here the main difficulty is already present in the
first-order statistic and standard error; using a larger Student critical
value only makes the conservative rejection rate smaller.

\begin{figure}[htbp]
\centering
\includegraphics[width=0.96\textwidth]{baseline_3x3_different_skew.pdf}
\caption[Baseline-MI comparison, \(3\times3\) tables]{Baseline-MI comparison for \(3\times3\) tables. Columns have
\(I(P)=0.0001,0.001,0.02\) nats, respectively; \(I(Q)=I(P)+|\Delta|\).
\(P\) has row and column margins \((0.7,0.15,0.15)\) and \(Q\) has
\((0.8,0.1,0.1)\). Both use the first-cell change block,
\(n_P=n_Q=100\), and \(|\Delta|=0,0.001,0.002,0.005,0.01,0.02\) nats.
Top: 0--1 rejection scale; bottom: 0--0.15 zoom. Each point has 20,000
replicates; hollow markers flag validity below 90\%.}
\label{fig:baseline-focus}
\end{figure}

Figure~\ref{fig:patterns-focus} asks whether the association being placed in
an early, rare or distributed set of cells changes performance when the
margins and target MI are fixed. At the displayed null points, Wald rates
are \(\PatternFirstWaldRate\), \(\PatternRareWaldRate\) and
\(\PatternSpreadWaldRate\) and Expanded Welch rates are
\(\PatternFirstExpandedRate\), \(\PatternRareExpandedRate\) and
\(\PatternSpreadExpandedRate\) for first, rare and spread patterns. The similarity of those three
null rates in this selected setting is not evidence that the patterns are
equivalent in all tables. The full \(3\times3\), \(5\times5\) and
\(8\times8\) matched sets are in the atlas.

Within this controlled comparison, moving the association to different cells
does not remove the conservative gap between the methods. The implication is
limited but useful: the displayed result is not an artefact of only the first
\(2\times2\) change block used in the main landscape.

\begin{figure}[htbp]
\centering
\includegraphics[width=0.96\textwidth]{patterns_5x5_different_skew.pdf}
\caption[Changed-cell comparison, \(5\times5\) tables]{Changed-cell comparison for \(5\times5\) tables. Columns place
association in the first \(2\times2\) block, last (rare) \(2\times2\)
block, or a centred score outer product spread across all cells. For every
column, \(P\)'s row and column margins are
\((0.7,0.075,0.075,0.075,0.075)\), and \(Q\)'s are
\((0.8,0.05,0.05,0.05,0.05)\). Here \(I(P)=0.0001\),
\(I(Q)=I(P)+|\Delta|\), \(n_P=n_Q=100\), and
\(|\Delta|=0,0.0001,0.00025,0.0005,0.001,0.002\) nats.
Top: 0--1; bottom: 0--0.15. Each point has 20,000 replicates.}
\label{fig:patterns-focus}
\end{figure}

\section{Unequal samples and the direction of the difference}

Figures~\ref{fig:imbalance-q} and~\ref{fig:imbalance-p} hold the unordered
sample sizes at 50 and 250 but reverse which population has the larger
sample. They also show the two directions of the MI difference separately.
The distinction matters because the finite bias corrections are
\(d/(2n_P)\) and \(d/(2n_Q)\), and the two populations have different
margins. For example, at \((n_P,n_Q)=(50,250)\) and a 0.02-nat difference,
Wald's rejection is \(\ImbalanceQHigherWaldRate\) when \(Q\) has higher MI and
\(\ImbalancePHigherWaldRate\) when \(P\) has higher MI\@. Expanded Welch gives
\(\ImbalanceQHigherExpandedRate\) and \(\ImbalancePHigherExpandedRate\). All four rates in
that specific comparison have validity one. These alternative rates should
still be read beside the matched null points; they are not equal-size power.

The practical implication is that an absolute MI difference does not define a
complete testing problem. Which population has the higher MI and which one has
the larger sample changes the bias and variance contributions. Power should
therefore not be averaged across the two directions of sample imbalance.

The larger-table atlas shows that this problem can be more severe than the
displayed \(3\times3\) examples. For the uniform \(8\times8\) null, changing
\((n_P,n_Q)\) from \((50,50)\) to \((50,500)\) increases the absolute
leading-correction difference from zero to
\(\ImbalanceEightCorrectionOffset\) nats. Wald rejection rises from
\(\ImbalanceEightEqualSmallWaldRate\) to
\(\ImbalanceEightQFiveHundredWaldRate\), and Expanded Welch rises from
\(\ImbalanceEightEqualSmallExpandedRate\) to
\(\ImbalanceEightQFiveHundredExpandedRate\); validity is one throughout. The
correction is intended to remove leading plug-in bias, so this pattern is
evidence that the first-order correction is under severe finite-sample stress,
not proof that the deterministic offset alone causes the inflation.

\begin{figure}[htbp]
\centering
\includegraphics[width=0.96\textwidth]{imbalance_3x3_q_higher.pdf}
\caption[Unequal samples, \(Q\) higher]{Unequal samples, \(Q\) higher. Columns have
\((n_P,n_Q)=(50,250)\) and \((250,50)\). For both columns, \(P\) has
\(3\times3\) row and column margins \((0.7,0.15,0.15)\) and \(Q\) has
\((0.8,0.1,0.1)\); the first-cell change block is used.
\(I(P)=0.02\) and \(I(Q)=0.02+|\Delta|\), with
\(|\Delta|=0,0.001,0.002,0.005,0.01,0.02\) nats. Top: 0--1; bottom:
0--0.15. Each point has 20,000 replicates.}
\label{fig:imbalance-q}
\end{figure}

\begin{figure}[htbp]
\centering
\includegraphics[width=0.96\textwidth]{imbalance_3x3_p_higher.pdf}
\caption[Unequal samples, \(P\) higher]{Unequal samples, \(P\) higher. Columns have
\((n_P,n_Q)=(50,250)\) and \((250,50)\). \(P\)'s \(3\times3\) row and
column margins are \((0.7,0.15,0.15)\); \(Q\)'s are \((0.8,0.1,0.1)\).
Both use the first-cell change block. \(I(Q)=0.02\) and
\(I(P)=0.02+|\Delta|\), with
\(|\Delta|=0,0.001,0.002,0.005,0.01,0.02\) nats.
Top: 0--1; bottom: 0--0.15. Each point has 20,000 replicates.}
\label{fig:imbalance-p}
\end{figure}

\section{Larger differences and rectangular tables}

The main grid deliberately concentrates on small MI differences. A separate
\(2\times2\) family extends the feasible axis to 0.2 nats without changing
the axis units or pretending unconstructed targets exist. This shows where
detection eventually rises, while retaining the zero-difference point that
defines null calibration. Rectangular \(2\times3\) and \(3\times5\) tables
check that the discussion is not tied only to square alphabets. The panels
do not average over shapes.

\begin{figure}[htbp]
\centering
\includegraphics[width=0.92\textwidth]{broad_2x2_uniform.pdf}
\caption[Wider MI differences, uniform \(2\times2\)]{Wider MI differences in a uniform \(2\times2\) population pair.
The single panel has \(P\) and \(Q\) row and column margins \((0.5,0.5)\),
first-cell changes, \(n_P=n_Q=100\), \(I(P)=0.02\), and
\(I(Q)=0.02+|\Delta|\). The horizontal grid is
\(|\Delta|=0,0.001,0.002,0.005,0.01,0.02,0.05,0.1,0.2\) nats;
the rejection axis is 0--1. Each point has 20,000 replicates.}
\label{fig:broad-focus}
\end{figure}

\FloatBarrier

The wider \(2\times2\) curve shows that both methods eventually detect large
MI differences, while Expanded Welch remains below Wald until the rejection
rates approach one. This is the expected cost of its larger critical value:
the power loss matters most at intermediate differences and becomes small when
the signal is overwhelming.

\begin{figure}[htbp]
\centering
\includegraphics[width=0.94\textwidth]{rectangular_different_skew.pdf}
\caption[Rectangular tables]{Rectangular tables. Columns are \(2\times3\) and \(3\times5\).
In each dimension \(k\), \(P\)'s corresponding margin is \(m_k(0.7)\)
and \(Q\)'s is \(m_k(0.8)\), for both rows and columns. Thus the first
column uses \(P\) rows \((0.7,0.3)\), columns
\((0.7,0.15,0.15)\), and \(Q\) rows \((0.8,0.2)\), columns
\((0.8,0.1,0.1)\); the second applies the same formula to 3 and 5
categories. Both use first-cell changes, \(n_P=n_Q=100\),
\(I(P)=0.02\), and \(I(Q)=0.02+|\Delta|\) for
\(|\Delta|=0,0.001,0.002,0.005,0.01,0.02\) nats.
Top: 0--1; bottom: 0--0.15. Each point has 20,000 replicates.}
\label{fig:rectangular-focus}
\end{figure}

\FloatBarrier

In the rectangular tables, both methods remain below 0.05 throughout the
displayed range, with Expanded Welch lower still. Those low alternative
rejection rates indicate weak detection, not good null calibration. Together
the two checks show that the main trade-off is not confined to square tables
or to the narrow MI grid.

\section{Population construction sensitivity}

The additive block is interpretable, but it is only one way to make tables
with fixed margins and target MI\@. Figure~\ref{fig:construction-focus}
compares it with an ordinal log-linear construction under matched targets.
The null rates are not identical: for the displayed \(5\times5\)
different-skew comparison at \(n=100\), Wald gives
\(\ConstructionAdditiveWaldRate\) under additive construction and
\(\ConstructionLoglinearWaldRate\) under log-linear construction; Expanded
Welch gives \(\ConstructionAdditiveExpandedRate\) and
\(\ConstructionLoglinearExpandedRate\). All are below nominal and all are
valid. Thus matching the scalar MI
and margins does not fix the complete finite-sample behaviour. The
constructor comparison is a sensitivity check, not proof that the population
families exhaust all relevant dependence structures.

\begin{figure}[htbp]
\centering
\includegraphics[width=0.94\textwidth]{construction_5x5_different_skew.pdf}
\caption[Additive versus log-linear construction]{Additive versus ordinal log-linear construction for
\(5\times5\) tables. The left column uses the first-cell additive block;
the right fits an ordinal interaction while preserving the margins. In both,
\(P\)'s row and column margins are
\((0.7,0.075,0.075,0.075,0.075)\) and \(Q\)'s are
\((0.8,0.05,0.05,0.05,0.05)\). Both have
\(I(P)=0.02\), \(I(Q)=0.02+|\Delta|\), \(n_P=n_Q=100\), and
\(|\Delta|=0,0.001,0.002,0.005,0.01,0.02\) nats.
Top: 0--1; bottom: 0--0.15. Each point has 20,000 replicates.}
\label{fig:construction-focus}
\end{figure}

\section{Extreme skew, independence and larger samples}

The extreme-skew check deliberately keeps configurations with very small
expected cell counts. At \(n_P=n_Q=100\), \(3\times3\) tables with dominant
margins 0.999 for \(P\) and 0.9995 for \(Q\) yield valid null results in
only 0.0118 of Wald replicates and 0.00005 of Expanded Welch replicates.
Both unconditional false-positive rates are zero, but this is mostly an
absence of valid tests, not evidence of excellent calibration. With dominant
margins 0.9 and 0.95, the corresponding null rates are 0.0701 and 0.0088;
Expanded Welch lowers an excessive Wald rate past the target.
Figure~\ref{fig:extreme-focus} shows all three specified skew levels, not just
the favourable one.

\begin{figure}[htbp]
\centering
\includegraphics[width=0.96\textwidth]{extreme_3x3.pdf}
\caption[Extreme-skew \(3\times3\) tables]{Extreme-skew \(3\times3\) tables. Across columns the dominant
row and column probabilities for \((P,Q)\) are
\((0.9,0.95)\), \((0.99,0.995)\), and \((0.999,0.9995)\). Each remaining
margin entry splits the leftover probability equally between the two other
categories. Both use first-cell changes, \(n_P=n_Q=100\),
\(I(P)=0.00001\), and \(I(Q)=I(P)+|\Delta|\), for
\(|\Delta|=0,0.00001,0.00005,0.0001,0.0002\) nats.
Top: 0--1; bottom: 0--0.15. Each point has 20,000 replicates; hollow
markers flag validity below 90\%.}
\label{fig:extreme-focus}
\end{figure}

Exact independence is a separate theoretical boundary.
Figure~\ref{fig:independence-focus} starts \(P\) at MI zero and includes the
equal-MI point where both populations are independent. This is not a routine
positive-MI validation of Expanded Welch: its population first-order
variance is zero at independence. At \(n=100\), for example, the uniform
\(2\times2\) null rates are only 0.00025 and 0.00010 for Wald and Expanded
Welch. Appendix~\ref{app:independence-boundary} explains why a second-order
independence test has a different reference.

\begin{figure}[htbp]
\centering
\includegraphics[width=0.96\textwidth]{independence_2x2.pdf}
\caption[Exact-independence boundary]{Exact-independence boundary for \(2\times2\) tables. Columns are
uniform margins \((0.5,0.5)\) for both populations, same skew
\((0.8,0.2)\) for both, and different skew with \(P=(0.7,0.3)\) and
\(Q=(0.8,0.2)\); each vector is used for both rows and columns. The
first-cell block is used, \(I(P)=0\), \(I(Q)=|\Delta|\),
\(n_P=n_Q=100\), and
\(|\Delta|=0,0.001,0.002,0.005,0.01,0.02\) nats.
Top: 0--1; bottom: 0--0.15. Each point has 20,000 replicates.}
\label{fig:independence-focus}
\end{figure}

Does increasing \(n\) resolve the null error when the populations remain
fixed? Figure~\ref{fig:convergence-focus} keeps equal MI and changes only
sample size. Across the 12 first-block profiles with \(I(P)=I(Q)=0.02\) and
\(n=50{,}000\), Wald's false-positive rates range from 0.0470 to 0.0527
and Expanded Welch's from 0.0466 to 0.0526. At equal MI 0.0001, however,
the ranges at the same \(n\) are 0.0089--0.0280 and 0.0076--0.0138.
These are finite-sample observations for fixed populations. They show that
very small positive MI can remain difficult at a large sample size, not that
the regular asymptotic theory is mathematically false.

\begin{figure}[htbp]
\centering
\includegraphics[width=0.9\textwidth]{convergence.pdf}
\caption[Large-sample null behaviour]{Large-sample null behaviour. Rows compare \(2\times2\) and
\(8\times8\) tables; columns have common true MI 0.0001 and 0.02 nats.
For both shapes, \(P\)'s row and column margins are \(m_k(0.7)\) and
\(Q\)'s are \(m_k(0.8)\), with \(k=2\) or 8. Both populations use the
first-cell block. The horizontal axis is the equal sample size
\(n_P=n_Q=1{,}000,2{,}500,10{,}000,50{,}000\) on a labelled log scale;
the vertical axis is null rejection from 0 to 0.10. Each point has 20,000
replicates.}
\label{fig:convergence-focus}
\end{figure}

\section{Mechanism checks and null composition}
\label{sec:mechanism-results}

The main-grid null contains two different hypotheses in terms of complete
population equality. The 72 uniform and same-skew configurations have
\(P=Q\), whereas the 36 different-skew configurations have \(P\ne Q\) but
\(I(P)=I(Q)\). The latter are the weak-null cases that most directly match
the thesis motivation. Table~\ref{tab:null-band-summary} reports these
profiles separately. The broad interval 0.025--0.075 is used only to describe
the landscape; it is not a formal acceptance region for calibration.

\begin{table}[htbp]
\centering
\small
\caption{Location, separately for each method, of the 108 main-grid null configurations relative to the descriptive rejection-rate interval 0.025--0.075. Each profile contains 36 configurations. The final column separately counts configurations with validity below 90\%; it overlaps the three rate-location columns. In particular, all four $n=2$ configurations per profile have zero validity for both methods. Uniform and same-skew are strong nulls; different-skew is the equal-MI weak null. Exact rates and validity are available in the saved results described in Appendix~\ref{app:evidence-index}.}
\label{tab:null-band-summary}
\begin{tabular}{@{}llrrrr@{}}
\toprule
Margins & Method & Below & Inside & Above & Validity below 90\%\\
\midrule
Uniform & Normal Wald & 10 & 18 & 8 & 4 \\
Uniform & Expanded Welch & 19 & 16 & 1 & 5 \\
Same skew & Normal Wald & 14 & 15 & 7 & 8 \\
Same skew & Expanded Welch & 25 & 11 & 0 & 12 \\
Different skew & Normal Wald & 9 & 16 & 11 & 8 \\
Different skew & Expanded Welch & 19 & 13 & 4 & 12 \\
\bottomrule
\end{tabular}
\end{table}


The table makes the direction of the adjustment clear. Expanded Welch greatly
reduces the number of rates above the descriptive band, but it also increases
the number below it in every margin profile. Its apparent protection against
liberal rejection is therefore purchased with additional conservatism rather
than with calibration in both directions.

A separate supplementary study tested the mechanism without altering the
confirmatory results. For the representative \(3\times3\) different-skew
null with \(I(P)=I(Q)=0.02\) and \(n=100\), Wald rejection was 0.0166. Using
an independent Monte Carlo estimate of the numerator's finite-sample standard
deviation gave 0.0551. In contrast, using the population first-order standard
deviation gave 0.0976. The estimated squared standard error was 1.17 times
the empirical variance, while the empirical variance was 1.42 times the
population first-order variance. Denominator estimation is therefore an
important source of conservatism in this setting, but the first-order
numerator approximation also matters.

The same follow-up compared the Hutcheson-style \(n_i\) assignment, Simple
Welch's \(n_i-1\) assignment and a kurtosis-only component degree of freedom
against the complete expansion. The first two were nearly indistinguishable
on average and generally close to Wald, although their rejection rates differed
by as much as 0.01235 in one \(5\times5\), \(n=5\) configuration. The
kurtosis-only method was less conservative than Expanded Welch,
showing that changes to the pointwise MI values and margins drive much of the
expanded adjustment. Appendix~\ref{app:mechanism} gives the derivation,
independent-pilot diagnostics, ablation figure and complete main-null table.

The convergence paths need not be monotone at the available finite sample
sizes. For the \(8\times8\), very-small-MI populations, the follow-up shows
that the ratio of estimated squared SE to empirical numerator variance changes
materially between \(n=2{,}500\), 10,000 and 50,000. As cells become occupied
and the quadratic interaction term gives way to the linear approximation,
the numerator distribution and denominator approximation can improve at
different rates. Asymptotic convergence does not require every intermediate
null rejection rate to move monotonically toward 0.05.

\section{Runtime and summary of evidence}

The complete scalar calls have the same \(O(rc)\) asymptotic complexity.
The observed median time across the 60 timing regimes is about 85
microseconds for Wald and 150 microseconds for Expanded Welch. The median
Expanded-to-Wald ratio is about 1.75, with regime ratios about 1.72--1.79.
For the explicitly timed \(3\times3\), different-skew, \(n_P=n_Q=100\),
equal-MI input regime, the medians are \(\RuntimeRegularWaldMedian\) and
\(\RuntimeRegularExpandedMedian\) microseconds.
These are implementation- and machine-specific timings, not universal
complexity constants. Exact per-regime medians and interquartile ranges are
in the saved runtime file cited in Appendix~\ref{app:evidence-index}.

Computationally, the correction is feasible: it has the same cell-linear order
as Wald and adds roughly 75\% to a scalar call in this implementation. The
overhead is modest compared with resampling, but it is still an extra cost.
Because the statistical results do not show a general improvement, runtime
does not provide a separate reason to prefer Expanded Welch.

\begin{table}[htbp]
\centering
\small
\caption{Complete scalar-call times in microseconds for the different-skew
null at \(n_P=n_Q=100\) and \(I(P)=I(Q)=0.02\). Entries give median
[interquartile range] over five passes on 200 saved inputs per regime.
Population construction and multinomial sampling are excluded.}
\label{tab:runtime-by-shape}
\begin{tabular}{@{}lccc@{}}
\toprule
Shape & Normal Wald & Expanded Welch & Expanded/Wald\\
\midrule
\(2\times2\) & 85.8 [85.5, 86.3] & 150.0 [149.4, 151.3] & 1.75\\
\(3\times3\) & 85.1 [84.6, 85.6] & 150.0 [148.9, 150.9] & 1.76\\
\(3\times5\) & 85.2 [84.9, 85.5] & 150.1 [149.5, 150.9] & 1.76\\
\(5\times5\) & 85.3 [85.1, 85.7] & 149.2 [148.5, 150.1] & 1.75\\
\(8\times8\) & 88.0 [87.1, 88.9] & 153.4 [152.4, 154.8] & 1.74\\
\bottomrule
\end{tabular}
\end{table}

The evidence separates three observations. First, Expanded Welch can make a
liberal Wald null rate closer to 0.05. Second, because it uses the same
statistic with a heavier-tailed reference, it can also reduce rejection where
Wald is already conservative. Third, in sparse samples the unconditional rate
can fall because the expanded calculation is invalid more often. No one
number summarises these behaviours; the matched curves and their validity
records are the basis for the discussion that follows.
